\section{Purgatory: Scripture and the Unbroken Confession of the Church}

Between immediate entrance into the beatific vision and immediate descent into eternal damnation lies a mercy belonging wholly to the saved. Many Christians die reconciled to God in sanctifying grace, united to Christ by faith and charity, yet still carrying unrepented venial faults, lingering attachments to created things, or an undischarged debt of temporal punishment left behind by forgiven sins. Because the Apocalypse states an absolute law for the heavenly Jerusalem---``There shall not enter into it any thing defiled or that worketh abomination or maketh a lie: but they that are written in the book of life of the Lamb'' (Apoc 21:27)---nothing stained can step directly into the vision of the all-holy God. Yet the soul dying in Christ's friendship cannot be lost in hell. The Church therefore confesses Purgatory: the final purification of the elect, by which grace cleanses every remaining blemish and brings redeemed souls into the unclouded joy of heaven.

Sacred Scripture grounds this faith across both Testaments. Its clearest Old Testament witness stands in the Second Book of Machabees. Discovering forbidden idolatrous tokens beneath the coats of fallen Jewish soldiers, Judas Machabeus turns immediately to prayer and Temple sacrifice:

\begin{quote}
And making a gathering, he sent twelve thousand drachms of silver to Jerusalem for sacrifice to be offered for the sins of the dead, thinking well and religiously concerning the resurrection. (For if he had not hoped that they that were slain should rise again, it would have seemed superfluous and vain to pray for the dead,) And because he considered that they who had fallen asleep with godliness, had great grace laid up for them. It is therefore a holy and wholesome thought to pray for the dead, that they may be loosed from sins. (2 Mach 12:43--46)
\end{quote}

Those slain had ``fallen asleep with godliness,'' yet still needed to be ``loosed from sins,'' and the living could aid them by prayer and expiatory sacrifice. From the beginning the Church recognized this passage as the biblical charter of her prayers for the faithful departed.\endnote{Cf.\ \textit{Catechism of the Catholic Church}, \S 1032, citing 2 Mach 12:46.}

Three New Testament passages converge with equal weight. First, our Lord warns that blasphemy against the Holy Ghost ``shall not be forgiven him neither in this world, nor in the world to come'' (Matt 12:32). St.\ Augustine points directly to what this phrasing implies: ``were there not some whose sins, though not remitted in this life, shall be remitted in that which is to come, it could not be truly said, `They shall not be forgiven, neither in this world, neither in that which is to come.'\,''\endnote{St.\ Augustine, \textit{The City of God} XXI.24, trans.\ Dods (1871).} St.\ Gregory the Great draws the same conclusion in his \textit{Dialogues}---cited by St.\ Thomas Aquinas and the \textit{Catechism of the Catholic Church}---teaching that certain lighter faults can be remitted after this life in the age to come.\endnote{St.\ Gregory the Great, \textit{Dialogues} IV.39, cited in Aquinas, \textit{ST} I-II, q.\ 87, a.\ 5, sed contra, and \textit{Catechism of the Catholic Church}, \S 1031.}

Second, in the Sermon on the Mount, Christ warns of a prison from which release comes only after full payment: ``lest perhaps the adversary deliver thee to the judge, and the judge deliver thee to the officer, and thou be cast into prison. Amen I say to thee, thou shalt not go out from thence till thou repay the last farthing'' (Matt 5:25--26). Unlike hell's everlasting pit, this prison has a terminus: \textit{donec reddas novissimum quadrantem}, until the last farthing is paid.

Third, St.\ Paul sets forth the image of the refiner's fire in his First Epistle to the Corinthians:

\begin{quote}
For other foundation no man can lay, but that which is laid: which is Christ Jesus. Now, if any man build upon this foundation, gold, silver, precious stones, wood, hay, stubble: Every man's work shall be manifest. For the day of the Lord shall declare it, because it shall be revealed in fire. And the fire shall try every man's work, of what sort it is. If any man's work abide, which he hath built thereupon, he shall receive a reward. If any man's work burn, he shall suffer loss: but he himself shall be saved, yet so as by fire. (1 Cor 3:11--15)
\end{quote}

The Apostle distinguishes three outcomes: he who destroys God's temple is destroyed (1 Cor 3:17); he who builds gold, silver, and precious stones receives an undiminished reward; and he who keeps Christ as his foundation yet builds wood, hay, and stubble suffers real loss (\textit{detrimentum patietur}), ``but he himself shall be saved, yet so as by fire''---\textit{ipse autem salvus erit, sic tamen quasi per ignem}. Malachias had already foretold the Messiah ``like a refining fire, and like the fuller's herb,'' sitting to purify and refine His servants ``as gold, and as silver'' (Mal 3:2--3).

The Church did not invent Purgatory in the Middle Ages; she prayed for the dead from the apostolic age, confessed post-mortem purification in the Profession of Faith drafted by Pope Clement IV in 1267 and sworn in the name of Emperor Michael VIII Palaeologus at the Second Council of Lyons (1274), and defined the doctrine solemnly in ecumenical councils at Florence and Trent. The confession at Lyons II affirmed that souls departing truly penitent in charity before completing worthy fruits of penance are cleansed after death by purifying punishments (\textit{poenis purgatoriis}), and that the suffrages of the living faithful---Masses, prayers, almsgiving, and other works of piety---avail for their relief.\endnote{Profession of Faith of Pope Clement IV (1267), sworn in the name of Emperor Michael VIII Palaeologus at the Second Council of Lyons (1274), DS 856.} At the Ecumenical Council of Florence (1439), the Decree of Union \textit{Laetentur caeli} solemnly defined those same two truths in common between East and West: post-mortem purification through purifying punishments (\textit{poenis purgatoriis}) and the real efficacy of the Eucharistic sacrifice and other suffrages offered by the living.\endnote{Council of Florence (1439), Decree of Union \textit{Laetentur caeli} (DS 1304).}

When sixteenth-century Reformers denied Purgatory and Masses for the dead, the Council of Trent defined Catholic faith across three sessions. In Session VI, Canon 30 (1547), the Council anathematized the claim that after receiving the grace of justification, guilt and eternal punishment are remitted in such a manner that no debt of temporal punishment remains to be discharged either in this world or in the next in Purgatory before entry into the kingdom of heaven.\endnote{Council of Trent, Session VI (1547), Decree on Justification, Canon 30 (DS 1580).} In Session XXII (1562), the Council defined that the Holy Sacrifice of the Mass is truly propitiatory and is rightly offered not only for the sins, punishments, satisfactions, and other necessities of the living faithful, but also for those who have died in Christ and are not yet fully purified.\endnote{Council of Trent, Session XXII (1562), Doctrine on the Sacrifice of the Mass, Chapter 2 and Canon 3 (DS 1743, 1753).}

In Session XXV (1563), the \textit{Decree on Purgatory} gathered this teaching into a single pastoral mandate. Appealing to Sacred Scripture, the ancient tradition of the Fathers, and prior councils, the Synod taught that Purgatory exists and that the souls detained there are assisted by the suffrages of the faithful and above all by the acceptable sacrifice of the altar. It instructed bishops to ensure that this sound doctrine be believed, maintained, and taught everywhere, while excluding from popular preaching subtle and difficult speculations that do not build up faith, and suppressing whatever smacks of vain curiosity, superstition, or base financial gain.\endnote{Council of Trent, Session XXV (1563), \textit{Decree on Purgatory} (DS 1820).}

The catechisms preserve that exact clarity. The \textit{Roman Catechism} speaks of ``the fire of purgatory, in which the souls of just men are cleansed by a temporary punishment, in order to be admitted into their eternal country, `into which nothing defiled entereth.'\,''\endnote{\textit{Catechism of the Council of Trent}, Part I, Creed Article V, trans.\ Donovan, p.\ 51.} The \textit{Catechism of the Catholic Church} reiterates the same dogma: those who die in God's grace and friendship while still imperfectly purified are already assured of their eternal salvation, yet undergo purification after death to attain the holiness required for heaven's joy---a final purification of the elect that the Church names Purgatory and distinguishes completely from the punishment of the damned.\endnote{\textit{Catechism of the Catholic Church}, \S\S 1030--1031.}

In the encyclical \textit{Spe salvi}, Pope Benedict XVI observes how deeply this doctrine corresponds to the lived condition of most believers: beneath repeated compromises with evil and the layers of sin accumulated in daily choices, an ultimate interior openness to truth, love, and God endures in the depths of the soul, so that the thirst for purity remains alive and asks for cleansing rather than destruction.\endnote{Pope Benedict XVI, Encyclical Letter \textit{Spe salvi} (30 November 2007), \S 46.}

Purgatory safeguards both God's holiness and His saving fidelity: He neither admits what is defiled into heaven nor abandons a fragile child who died clinging to His Son.
